\section{Replacing the density by unit nuclei}\label{sec:nuclei}

Our input is the nonnegative density in~\eqref{eq:4}. We must replace its
Coulomb field by finitely many equal masses $1/\lambda$, so that physical
nuclear charges are all one. The construction has two distinct error
requirements: a form bound on arbitrary states and a stronger
compression bound on the smooth localized modes.

The comparison with \(H_d\) uses energies multiplied by \(\lambda\).
A compressed error must therefore be smaller than \(\lambda^{-1}\),
up to the required polynomial margin. The error coupling the mode space
to its complement can be larger: the fixed complementary gap makes its
energy contribution quadratic, as proved in
Section~\ref{subsec:complement-control}. With
\(h=(8/(\rho\lambda))^{1/3}\), a cubature rule exact through degree
three will give an \(O(h^4)\) compressed error and an \(O(h^2)\) form
error. Both resulting energy errors are \(O(\lambda^{-1/3})\) in
the units of \(H_d\), up to polynomial factors.

\subsection{Transport and charge placement}\label{subsec:charge-placement}

The density interpolation uses Moser's deformation
method~\cite[Section 4]{Moser1965}, also developed by Dacorogna and
Moser~\cite[Section II, Lemma 3]{DacorognaMoser1990}. Here the vector
field is explicit and compactly supported inside the box; the following
continuity-equation calculation proves the required transport directly.
Write
\(q_t=\rho+t\Delta V_\ell/(4\pi)\), so that \(q_0=\rho\)
and \(q_1\) is the desired density on \(\Omega\). Since
\(q_t\ge\rho/2\), the time-dependent vector field on full space
\[
 \mathcal V_t(b)=-\frac{\nabla V_\ell(b)}{4\pi q_t(b)},
 \qquad 0\le t\le1,
\]
has uniformly bounded derivatives through order four. It vanishes
outside the well cylinders strictly inside the box. Let \(G_t\) be its
flow from time zero and put \(G=G_1\). The ODE and its differentiated
equations give bounded derivatives of \(G\) through order four and a
bounded Lipschitz inverse. The flow is the identity off the cylinders
and preserves \(\Omega\).

The identity \(\partial_t q_t+\nabla\cdot(q_t\mathcal V_t)=0\)
gives
\[
 q_t(G_t(b))\det DG_t(b)=\rho,\qquad
 \int_\Omega F(b)q_1(b)\,db
       =\rho\int_\Omega F(G(b))\,db.
\]
Thus the flow transports the uniform measure to the target density.
The last formula also shows how to retain equal cubature masses:
move the nodes by \(G\) and keep their weights unchanged.

Since \(H,S,k\) are integers, cubes of side \(h=1/k\) tile
\(\Omega\). On each cube use the eight tensor products of the two
equal-weight Gauss points per coordinate, at the midpoint plus or minus
\(h/(2\sqrt3)\). The rule integrates polynomials of total degree at
most three exactly, by the corresponding moments on each interval.
Map these nodes by \(G\). Their common mass remains
\(\rho h^3/8=1/\lambda\); polynomial exactness is used in the
original cube coordinates, on integrands composed with \(G\).

Each exact mapped point \(b_\alpha\) will be approximated within
Euclidean error at most
\[
                    \delta_R\le (1+H)^{-3}(1+S)^{-2}\lambda^{-3}
\]
by rational points \(\tilde b_\alpha\), and the physical nuclear positions are \(\tilde b_\alpha/\lambda\). They remain distinct by node separation and bi-Lipschitz bounds.

\subsection{Form and compression estimates}\label{subsec:nuclear-errors}

Let \(L_p\) be the shifted one-particle operator with these actual point nuclei in dilated units, i.e. with potential \(-\sum_\alpha \lambda^{-1}/|y-\tilde b_\alpha|-C_H-e_*\). \begin{proposition}[Unit-nucleus replacement]\label{prop:nuclear-replacement}
For the transported degree-three cubature and rational rounding just
specified, with $\delta_R\le(1+H)^{-3}(1+S)^{-2}\lambda^{-3}$, the
shifted continuous and point-nucleus operators $L_c,L_p$ satisfy
\[
 \begin{split}
 |\langle u,(L_p-L_c)u\rangle|
      &\le p(n,D)\lambda^{-2/3}(\|u\|^2+\|\nabla u\|^2),\\
 \|P(L_p-L_c)P\|_{P}
      &\le p(n,D)\lambda^{-4/3}.
 \end{split}                                               \tag{7}\label{eq:7}
\]
The first estimate holds for every one-particle $H^1$ function $u$.
The second is the norm of the compressed form on the localized-mode
space $\operatorname{Ran}P$. Both bounds are in dilated energy units.
\end{proposition}

\begin{proof}
We first work at the exact transported nodes. On an original grid cube
\(Q\), write \(a_{Q,\nu}\), \(1\le\nu\le8\), for its Gauss
nodes. Taylor's theorem and exactness through degree three give the
common error estimate
\[
 \left|\rho\int_Q F(G(b))\,db
   -\frac{\rho h^3}{8}\sum_{\nu=1}^8 F(G(a_{Q,\nu}))\right|
 \le C h^7\sup_Q |D^4(F\circ G)|
\]
whenever \(F\circ G\) has four bounded derivatives on \(Q\).
We apply this first to the singular Coulomb kernel away from its
singularity, then to the smooth fields obtained after compression.

\noindent\textbf{General form estimate.} For fixed evaluation point \(y\), the charge integral is of the kernel \(|y-G(b)|^{-1}\) against uniform density in \(b\). For cubes whose images touch a ball of radius a fixed large constant times \(h\) about \(y\), there are only \(O(1)\) such cubes. Their continuous contribution is \(O(h^2)\) by bounded transformed density. The discrete one is bounded by a sum of \(C h^3/|y-b_\alpha|\) with \(|y-b_\alpha|\le C'h\). 

On farther cubes, the pure fourth derivative of the kernel is bounded
by \(C|y-G(b)|^{-5}\). Distances to \(y\) are comparable within
each such cube, and the bi-Lipschitz change of variables has bounded
Jacobian. Summing the pure derivative contributions therefore costs
at most
\[
 C h^4\int_{|x-y|>ch}|x-y|^{-5}\,dx=O(h^2).
\]
The other chain-rule terms contain a derivative of \(G\) of order
at least two and a kernel derivative bounded by
\(C|y-G(b)|^{-j}\), \(2\le j\le4\). Such terms occur only on
cells meeting the well cylinders, of combined volume
\(O(m(S+1))\). Below distance one they are bounded by the same
inverse fifth power; beyond distance one their contribution is at
most \(C h^4m(S+1)\).

This proves a pointwise absolute error bound by \(C(h^2+h^4 m(S+1))\) plus the close-node sum above. In particular there is no factor the size of the large horizontal area in this bound.

The close-node balls have bounded overlap also upon any fixed dilation. Around each node one has
\[
  \int_{|y-b_\alpha|\le C'h}\frac{|u(y)|^2}{|y-b_\alpha|}\,dy
  \le C\int_{|y-b_\alpha|\le 2C'h}
                       (h^{-1}|u|^2+h|\nabla u|^2).
\]
This follows by a cutoff and the three-dimensional Hardy inequality \(\int |u|^2/|y|^2\le4\int|\nabla u|^2\), which follows also by the divergence formula for \(y/|y|^2\) and Cauchy-Schwarz. Summing gives the general form bound, since \(h^2=O(\lambda^{-2/3})\).

\noindent\textbf{Compressed estimate.} For compression use the smooth Coulomb fields
\[
     K_{ij}(b)=\int |y-b|^{-1}\psi_i(y)\psi_j(y)\,dy .
\]
Their derivatives through order four are uniformly bounded (differentiate the translated densities). Fourth derivatives satisfy
\[
         |D_b^4 K_{ij}(b)|\le C(1+|b-(u_i,0)|)^{-5}.
\]
Indeed derivatives of the product density decay exponentially about \((u_i,0)\), uniformly by bounded derivatives on the other factor. At a large distance \(d=|b-(u_i,0)|\), use a smooth density cutoff supported within distance \(d/2\) of that center and equal to one inside \(d/3\), with the cutoff radius held fixed when differentiating. There one uses pure kernel derivative decay, and on the remaining tail one differentiates the cut-off density instead, with exponentially small bounds. This proves the assertion. Applying the same degree-three cubature to \(K_{ij}\circ G\) now costs \(O(h^4(1+m(S+1)))\) for each entry: pure fourth derivatives have summable volume bounds as just proved, and other chain-rule terms live on the exceptional cells. Gram correction and taking matrix norms cost at most polynomial factors. This proves the second estimate for exact nodes.

\noindent\textbf{Rounding.} A center shift of size \(\delta_R\)
changes the inverse distance by at most
\(\delta_R/(|y-b_\alpha|\,|y-\tilde b_\alpha|)\).
Cauchy--Schwarz and translated Hardy inequalities bound its form
by \(4\delta_R\|\nabla u\|^2\). Summing the masses gives the
total bound \(4\rho|\Omega|\delta_R\|\nabla u\|^2\), where
\[
 \rho|\Omega|\delta_R
 \le C H^{-1}S^{-1}\lambda^{-3}
 \le C\lambda^{-8}.
\]
This is smaller than both required errors, also after compression
because the modes have polynomial kinetic bounds. It completes
\eqref{eq:7}.
\end{proof}

We have obtained the two replacement estimates with constants
independent of the horizontal area of the slab. The full mass enters
only the rounding bound, where the chosen coordinate accuracy absorbs
it. Section~\ref{sec:many-body} now combines these estimates with the
fixed complementary gap and the compressed matrix from
Section~\ref{sec:density}.
